\section{Method}
\label{sec:method}

% ---- 草稿说明（改完删除）----------------------------------------------------
% Claude 起草的粗稿，供修改。PDF 中蓝色 [Claude: ...] 是需要你决定或核对的地方，
% \draftfalse 后自动消失（\note 内只能写英文，pdfLaTeX 不支持中文）。
% 目前只有 4.1--4.3，不含 Setup 段，也不引用尚未写的 Sec. 3 / Sec. 5 内容。
% 删掉的 Setup 段与原 4.4（训练目标、straight-through、两阶段训练、fold、计算量公式）
% 可在 git 提交 9935f2e 中找回。
% 所有超参数取自已采用格（contiguous-t，冻结 v2-gtmix router + 新训 LoRA 1500 步）。
% ---------------------------------------------------------------------------

\begin{figure*}[t]
  \centering
  \placeholder{4.5cm}
  % \includegraphics[width=\linewidth]{overview}
  \caption{\textbf{Overview of \method.}
  After an early block of the frozen video DiT, an action-aware router scores the future tokens.
  Spatially, the highest-scoring $2{\times}2$ blocks keep full resolution and every other block is merged into a single tube;
  temporally, each tube is pooled into a number of segments set by its relevance, and the segments are re-indexed contiguously in time.
  The remaining blocks run on the compact sequence, whose output conditions the action decoder.
  \TODO{overview figure}}
  \label{fig:overview}
\end{figure*}

Our method aims at identifying action-relevant tokens throughout within future representation, which then merge the other tokens to eliminate its negative influences towards both action generation and inference time. Specifically, our method compresses the future-frame tokens of the video backbone after an early block.
It decides, per region of the scene, at which spatial and temporal resolution the backbone computes, using relevance learned from the action decoder.

% ---------------------------------------------------------------------------
\subsection{Action-Aware Router}
\label{sec:router}

The router is a lightweight module that scores how relevant each future token is to the action.
It takes the future-frame tokens after an early block of the video DiT, together with the instruction embedding,
and outputs a score between 0 and 1 for every token:
each token attends to the instruction through one cross-attention layer, and an MLP with a sigmoid produces the score.

Since the allocation is decided once for the whole prediction horizon, we summarize the scores per spatial region.
We take the maximum score of each position over the future frames, and average it over each $2{\times}2$ block of the token grid:
\begin{equation}
  b_j = \operatorname*{mean}_{p \in \text{block } j} \; \max_t \, g_t(p),
  \label{eq:block_score}
\end{equation}
where $g_t(p)$ is the score of the token at position $p$ in future frame $t$.

The router does not measure visual similarity.
It is action-aware because it is trained from the action decoder:
its supervision is how much each block affects the decoder's output.
\note{Claude}{Do not call the router instruction-dependent or dynamic: on our diagnostics swapping the instruction leaves the allocation almost unchanged (Spearman 0.96), i.e.\ it behaves as a learned, near-static prior per suite.}
\note{Claude}{The router architecture and its pseudo-label training recipe come from an existing KnowledgeRouter implementation; cite its source if it is published work.}

% ---------------------------------------------------------------------------
\subsection{Spatial Resolution Allocation}
\label{sec:spatial}

We divide the $30{\times}40$ token grid into 300 blocks of $2{\times}2$ tokens and rank them by the block score $b_j$.
The 100 highest-scoring blocks keep all four tokens.
Every other block is merged into a single token by averaging its four tokens, separately in each frame.
This leaves 600 spatial positions, each followed through the 14 future frames; we call such a sequence a \emph{tube}.

No region is removed: relevant regions keep full resolution and the rest are represented at a coarser one.

\paragraph{Spatial positions.}
Every token keeps its position on the original grid.
A merged token takes the position of the top-left token of its block.
We do not re-index the tokens onto a smaller grid, so the backbone still sees where each region lies in the scene.

% ---------------------------------------------------------------------------
\subsection{Temporal Resolution Allocation}
\label{sec:temporal}

Each tube is then shortened along time.
We split its 14 frames into a few consecutive segments and average the tokens within each segment, so a tube with $K$ segments contributes $K$ tokens.

Tubes receive different numbers of segments according to their relevance.
We rank the tubes by the score of their block, weighted by how much the tube's features change between consecutive frames.
The most relevant tubes receive the most segments, using a fixed number of tubes per level.
For example, when 10\% of the future tokens are kept, the top 300 tubes receive 4 segments, the next 180 receive 2, and the last 120 receive 1, giving 1,680 tokens.
Because the numbers per level are fixed, every sample has exactly the same token count.

Within a tube, segment boundaries are placed so that each segment covers a similar amount of feature change.
% 按现行实现写：排序用 router 分数 x 帧间特征变化（cosine 距离），切点按累计变化的等分位（动态规划求解）。
% 按约定不写"按运动自适应"、不声称优于固定 K。
% 诊断事实：K 档与 router 的 fine/coarse 划分完全重合，给定 K 时切点近乎固定（K=4 时段长 4/3/3/4），
% 部署噪声固定（seed 0）时不同观测间 K 分配 97.3% 相同。即分配实际上由 router 决定。
\note{Claude}{Written as implemented. In practice the levels coincide with the router's full/merged split and the cuts are near-uniform, so the allocation is effectively decided by the router; decide whether to drop the feature-change weighting from the text and describe the cuts as even splits.}

\paragraph{Temporal positions.}
Unlike the spatial axis, temporal positions are re-indexed:
the $k$-th segment of every tube takes the $k$-th future time position, whichever frames it covers.
Temporal positions therefore count segments rather than original frames, the opposite of the spatial convention.
